% Separate post-CP20 candidate. Not part of any earlier checkpoint.
\subsection{A cubic wild layer and the residual projective order}

\begin{lemma}[The anti-Cartan layer and exact tame coordinates]
\label{lem:B-RD56-anti-layer}
\status{Candidate; complete proof written, independent audits pending.}
Use the normalized remaining $p^2$ comparison of
Theorem~\ref{thm:B-RD49}, with $p\ge5$, over a complete discretely
valued mixed-characteristic $(0,p)$ field $K$ with algebraically
closed residue field, $e_K=p-1$, and $\mu_p\subset K$.
Assume that the common residual representation is irreducible over
$\mathbf F_p$, that its projective image has order $h>2$, and that
the two special characters differ. Put $q=p^2$ and $N=q-1$.
If its common second-level wild image $W$ has order $p^3$, then
$W$ contains a distinguished two-dimensional subspace $U$ of
anti-Cartan matrices. Conjugation by tame inertia acts irreducibly
on $U$. All its nonzero elements have the same lower break $b$.

Take the common tame good field $L/K$ of degree $N$ and the common
second-level division field $F/L$. One can choose a tame complement
$\langle\sigma\rangle\subset\operatorname{Gal}(F/K)$, a uniformizer
$\varpi$ of $F$, and formal coordinates $T_i$ for the two good
models such that, for a primitive $N$th root $\zeta$ and their
actual tangent exponents $c_i$,
\[
 \begin{gathered}
 \sigma\varpi=\eta\varpi,\qquad \eta=\zeta^p,\qquad
 T_i(\sigma P)=\zeta^{c_i}\sigma(T_i(P)),\\
 A=v_F(T_i(P))=p(p-1)\quad(P\text{ primitive of order }p^2),
 \qquad v_F(T_i(Q))=qA\quad(0\ne Q\in M_i[p]),\\
                         e_F=NqA.                       \tag{B-RD56.1}
 \end{gathered}
\]
The special coordinates remain Honda coordinates with $[p](T)=T^q$.
\end{lemma}
\begin{proof}
Identify either residual space with $\mathbf F_q$, so its tame
image acts by multiplication by $\lambda\in\mathbf F_q^\times$.
Every endomorphism is uniquely
\[
                         z\longmapsto az+bz^p.
\]
Conjugation fixes the two-dimensional Cartan space of $a$'s and
acts on the two-dimensional anti-Cartan space of $b$'s by
$\lambda^{1-p}$. This element has order $h>2$; since its order
divides $p+1$, it is not in $\mathbf F_p$. Thus the anti-Cartan
space is irreducible over $\mathbf F_p$. The conjugation action
has prime-to-$p$ order and is semisimple by averaging. Its invariant
three-dimensional subspace $W$ must consequently be the direct
sum of this anti-Cartan space $U$ and a Cartan line.
The lower ramification filtration is invariant under conjugation.
Its intersections with the irreducible space $U$ are either all
of $U$ or zero, giving a single break $b$ for its nonzero elements.

Proposition~\ref{prop:B-RD53-cycle} gives the common tame good
field of degree $N$. The argument of Lemma~\ref{lem:B-RD55-valuations}
gives $\operatorname{Gal}(F/L)=W$ and the displayed valuations.
The extension of $C_N$ by the abelian group $W$ has a complement:
if a lift $s$ of a generator has $s^N=w\in W$, then $w$ is fixed
by $s$, and replacing $s$ by $s(-N^{-1}w)$ makes its $N$th power
one. Here $W$ is written additively and $N$ is invertible modulo $p$.
Choose the generator $\sigma$ to act on a tame uniformizer of $L$
by $\zeta$. Its multiplier $\eta$ on a uniformizer of $F$ satisfies
$\eta^{p^3}=\zeta$, hence $\eta=\zeta^p$. Averaging a uniformizer
against this character makes $\sigma\varpi=\eta\varpi$ exact,
without changing its leading term.

The actual semilinear descent isomorphisms extend to the good models
by Tate full faithfulness, as in Lemma~\ref{lem:B-RD53-leading}.
Choose the special Honda framings in which the finite tame actions
are $T\mapsto\zeta^{c_i}T$. Averaging any integral lift of that
coordinate against this character over the cyclic group of order
$N$ gives the exact semilinear eigen-coordinate in (B-RD56.1).
Its special reduction is unchanged, so its special $p$-series is
still $T^q$. This uses the actual descent cocycle; it is not an
assumption that an arbitrary deformation has CM.
\end{proof}

\begin{proposition}[Exclusion of order $p^3$ when the projective order exceeds two]
\label{prop:B-RD56-cubic-wild}
\status{Candidate; complete proof written, independent audits pending.}
Under the normalized hypotheses of the preceding lemma, if the
common residual projective image has order greater than two,
isomorphic $p^2$-torsion with different special characters must
have wild image of order $p^4$. In particular wild order $p^3$
is excluded, in addition to the order $p^2$ already excluded by
Proposition~\ref{prop:B-RD55-minimal-wild}.
The same exclusion holds when the residual projective order is two,
provided that $W$ contains the entire anti-Cartan plane and all its
nonzero elements have one common lower break.
\end{proposition}
\begin{proof}
Suppose instead that $|W|=p^3$ and first use the preceding lemma.
For the additional projective-order-two assertion, take the full
anti-Cartan plane $U$ assumed in the statement. The same common field,
valuation and exact-coordinate construction applies: only the lemma's
automatic existence and single-break assertion used $h>2$.
Fix a coefficient ring in $\mathcal O_K$ and residue representatives
as in (B-RD55.4). Since $e_F=NqA>qA$, reduction modulo
$\varpi^{qA+1}$ identifies these coefficients with the field $k$:
all errors in their addition and multiplication are divisible by
$p$ and vanish in this quotient. We may therefore calculate with
characteristic-$p$ power series to this precision.

Let $R_i=M_i[p^2](F)$ and $V_i=R_i/pR_i$. Primitive leading
coefficients give $\ell_i:V_i\hookrightarrow k$ with image
$d_i\mathbf F_q$; multiplication by $p$ raises leading coefficients
to the $q$th power, exactly as in the proof of
Proposition~\ref{prop:B-RD55-minimal-wild}. Write $\ell_i(v)=d_i z$.
Changing a representative of $v$ adds a first-level point, so it
changes its formal coordinate only in valuation at least $qA$.
Thus, for all $j<qA$, the coefficients in
\[
 T_i(P)=d_i z\varpi^A+
                \sum_{A<j<qA} f_{i,j}(z)\varpi^j
                     \pmod{\varpi^{qA}}                \tag{B-RD56.2}
\]
are well-defined functions of $z\in\mathbf F_q$.
No additivity of every coefficient function is assumed.

For $g\in U\setminus\{0\}$, put
\[
 d_g=\overline{(g\varpi-\varpi)/\varpi^{b+1}}.
\]
Composing two such expansions adds their leading coefficients;
the kernel is $U\cap G_{b+1}=0$. Thus $g\mapsto d_g$ is
injective and additive on $U$. Applying
$\sigma g\sigma^{-1}$ to $\varpi$ shows directly that tame
conjugation multiplies the leading coefficient by $\eta^b$.
This realizes the two-dimensional conjugation module $U$ inside $k$.
When $h>2$ that module is irreducible; when $h=2$ conjugation acts
on the entire anti-Cartan plane by the scalar $-1$. In either case
its action is nontrivial, so
\[
                         \eta^b\ne1.                  \tag{B-RD56.3}
\]
Put $b_0=N(p-1)$. Equation (B-RD56.3) implies $b\ne b_0$.
We use $b_0$ only as this numerical comparison value; no identification
of a lower break with a cyclotomic upper break is made.

Every nonzero anti-Cartan matrix $z\mapsto\beta z^p$ is invertible.
Consequently $gP-P$ is a nonzero first-level point for every primitive
$P$. Formal subtraction differs from ordinary coordinate subtraction
by a unit with residue one, so
\[
                        v_F(gT_i(P)-T_i(P))=qA.          \tag{B-RD56.4}
\]
For a monomial of degree $j$, if $p^s$ is the exact power of $p$
dividing $j$, the first term of its change has degree
$j+p^s b$ and coefficient $(j/p^s)d_g^{p^s}$ modulo $p$.
This follows by writing $(\varpi+d_g\varpi^{b+1}+
\text{higher terms})^j$ in characteristic $p$.
Terms divisible by $p$ in mixed characteristic have valuation
strictly larger than $qA$ and cannot affect this calculation.

Suppose first that $b<b_0$. The degree-$A$ monomial in (B-RD56.2)
first changes in degree $A+pb<qA$, with nonzero coefficient
$(p-1)d_i z d_g^p$. Every other monomial whose degree is divisible
by $p$ first changes strictly later. Among monomials of degree
prime to $p$, take the first with a nonzero coefficient. If its
degree were less than
\[
                           j_0=A+(p-1)b,
\]
its change would be the unique earlier nonzero term, contradicting
(B-RD56.4). A later degree cannot cancel the degree-$A$ term.
Hence $p\nmid j_0$, and cancellation forces, for every $z$,
\[
             j_0 f_{i,j_0}(z)
                     =-(p-1)d_i d_g^{p-1}z.            \tag{B-RD56.5}
\]
Fix any $z\ne0$. The left side of (B-RD56.5) is independent of
$g$, whereas $d_i z\ne0$. Hence $d_g^{p-1}$ would have to be the
same nonzero constant for all $g\in U\setminus\{0\}$.
The injective map $g\mapsto d_g$ has $q-1>p-1$ distinct nonzero
values, but a polynomial $X^{p-1}-c$ has at most $p-1$ roots.
This contradiction excludes $b<b_0$. It needs neither a CM assumption
nor a choice of a one-dimensional character inside $U$.

Suppose now that $b>b_0$. Every monomial of degree divisible
by $p$ changes in degree at least $A+pb>qA$.
Equation (B-RD56.4) therefore forces the first nonzero coefficient
of prime-to-$p$ degree to have
\[
                         j=qA-b,\qquad p\nmid j.
\]
Here $A<j<qA$ is necessary; if this fails, (B-RD56.4) is
already impossible. Its leading displacement is
$j f_{i,j}(z)d_g$.
Write the matrix of $g$ on $R_i$ as $1+pB_i(g)$.
On $U$ it has $B_i(g)(z)=\beta_i(g)z^p$, and
$\beta_i:U\to\mathbf F_q$ is an additive isomorphism.
The first-level leading coefficient comparison gives
\[
               d_i^q\beta_i(g)z^p
                           =j f_{i,j}(z)d_g.            \tag{B-RD56.6}
\]
Fixing one nonzero $g$ shows that $f_{i,j}(z)$ is a fixed
nonzero multiple of $z^p$. Applying the equation to every
$g\in U$ then gives constants $\kappa_i\ne0$ with
\[
                          \beta_i(g)=\kappa_i d_g.
                                                        \tag{B-RD56.7}
\]

Because the special characters differ, the common residual
intertwiner in these leading coordinates is $z\mapsto c z^p$,
by Lemma~\ref{lem:B-RD53-leading}. Multiplication by $p$
identifies the induced maps on $V_i$ and $pR_i$, since the
normalized leading-coordinate map is $z\mapsto z^q=z$.
Conjugation of the anti-Cartan matrices consequently gives
\[
                      \beta_2(g)=c^{1-p}\beta_1(g)^p.
\]
Together with (B-RD56.7), this makes Frobenius a fixed scalar
multiplication on $\beta_1(U)=\mathbf F_q$, an impossibility.
This excludes $b>b_0$ and proves the proposition.
\end{proof}

\begin{corollary}[A smaller wild list at five]
\label{cor:B-RD56-p5}
\status{Candidate; independent audits pending.}
In the $p=5,e_K=4$ remaining test with irreducible residual
inertia, the pattern $\{A,B\}$ of Corollary~\ref{cor:B-RD51-p5}
requires wild order $625$. The pattern $\{A,C\}$ still has
only the necessary wild list $125$ or $625$ from B55.
Neither remaining case is asserted to exist.
\end{corollary}
\begin{proof}
For $\{A,B\}$, the actual tangent exponents giving equal
residual characters are $1,21$, up to interchange and simultaneous quadratic twist. The residual eigenvalues have order $24$
and their projective ratio has order $6>2$.
For $\{A,C\}$, the equal-residual tangent exponents are $5,17$
or their simultaneous variants; the residual eigenvalues have
order $8$ and projective ratio order $2$.
Apply the proposition in the first case and retain the exact
scope of Proposition~\ref{prop:B-RD55-minimal-wild} in the second.
\end{proof}

\begin{remark}[Scope and numbering caution]
\label{rem:B-RD56-scope}
\status{Candidate scope statement.}
The argument concerns the actual irreducible residual representation
and the common second-level wild Galois module. It does not assume
that the groups have CM or that one is a full Tate-module twist of
the other. Reducible residual representations and wild order $p^4$ are not
excluded here. At projective order two, the new assertion excludes
only an entire anti-Cartan plane with a single common lower break;
a one-dimensional anti-Cartan part or two distinct anti-Cartan
breaks are not settled by this statement. The intermediate outline's
suggestion to identify a Cartan lower break with a cyclotomic upper
break is not used: the proof works entirely with lower breaks and the
explicit comparison number $b_0=(p^2-1)(p-1)$.
The exact finite supplement
\path{code/reduction_descent/verify_tame_cubic_wild_layer.py}
checks $2{,}498$ monomial first terms, $2{,}494$ lower-break cases,
$2{,}336$ tame-weight cases, and $45{,}864$ explicit anti-Cartan
conjugacy cases. It reproduces all six oriented normalized cases
at $p=5$. The original author run passed in $0.4074130669978331$ seconds.
Its source and output are retained before the additional low-break
cardinality checks. The expanded run passed in
$0.4129977519987733$ seconds, additionally checking the nonconstancy
of $d^{p-1}$ on $360$ nonzero elements in four quadratic finite fields.
This is finite arithmetic support for the proof, not a formal-group
construction or a pair-existence computation.
\end{remark}
